\[
|\hat{p}_{i1}-0.683|>3\times\sqrt{\frac{0.683\times0.317}{n}}=\frac{1.396}{\sqrt{n}}
\]
或者
\phantomsection\label{eq:4-29}
\begin{equation}
|\hat{p}_{i2}-0.954|>3\times\sqrt{\frac{0.954\times0.046}{n}}=\frac{0.628}{\sqrt{n}}
\tag{4-29}
\end{equation}
表明了对第 $i$ 个特性的正态分布假定的偏离. 当观测比例太小时，建议采用比正态粗尾的母体分布.

在任何数据分析中，画图总是有效的方法. 有一类可专门用于评估正态性假定的图形，叫做 $Q$-$Q$ 图. 对每个变量的样本观测值的边缘分布可以作此图. 实际上，它们展示的是样本分位数与观测值确实是正态分布时所会观察到的分位数之间的关系. 当各点离一条直线很近时，正态性假定仍然是保持的；若点偏离直线，正态性就较为可疑. 而且，偏差的图形可提供有关非正态性种类的线索. 一旦非正态的原因被确定，修正工作通常是可能的（参见 4.8 节）.

为了简化记号，令 $x_{i1},x_{i2},\cdots,x_{in}$ 表示任何单一特征 $X_i$ 的 $n$ 个观测值. 这些观测值按大小排序以后表示为 $x_{i(1)}\leqslant x_{i(2)}\leqslant\cdots\leqslant x_{i(n)}$. 例如，$x_{i(2)}$ 是第二小的观测值，$x_{i(n)}$ 是最大的观测值. 当 $x_{ij}$ 互不相同时，恰好有 $j$ 个观测值小于或等于 $x_{i(j)}$，$x_{i(j)}$ 是样本分位数（当观测值是我们通常假定的连续型时，这在理论上总是成立的）. 为了分析方便，样本在位于 $x_{i(j)}$ 的左边的比例 $j/n$ 通常用 $\left(j-\dfrac{1}{2}\right)\big/n$ 近似.\footnote{$(j-1/2)/n$ 的分子中，$1/2$ 是一个“连续性”修正. 有些作者（参见文献[5]和文献[10]）建议用 $(j-3/8)/(n+1/4)$ 代替 $(j-1/2)/n$.}

对于标准正态分布，分位数 $q_{i(j)}$ 由下述关系定义
\phantomsection\label{eq:4-30}
\begin{equation}
P[Z\leqslant q_{i(j)}]=\int_{-\infty}^{q_{i(j)}}\frac{1}{\sqrt{2\pi}}e^{-z^{2}/2}dz=p_{i(j)}=\frac{j-\frac{1}{2}}{n}
\tag{4-30}
\end{equation}
（参见附录中的表 1.）这里 $p_{i(j)}$ 是从标准正态总体取一次值会小于或等于 $q_{i(j)}$ 的概率.

基本想法是在有关的累积概率 $(j-1/2)/n$ 相同的情况下，考察分位数对 $(q_{i(j)},x_{i(j)})$. 若数据由正态总体产生，由于 $\sigma q_{i(j)}+\mu$ 近乎是期望的样本分位数，点对 $(q_{i(j)},x_{i(j)})$ 会近似线性相关.\footnote{一个较好的方法是画出 $(m_{i(j)},x_{i(j)})$，其中 $m_{i(j)}=E(x_{i(j)})$ 是从一个标准正态分布中抽取容量为 $n$ 的样本的第 $j$ 个顺序统计量的期望值（参见文献[12]以便进一步讨论）.}

\noindent\textbf{例 4.9}\quad (作 $Q$-$Q$ 图)

由下表给出一个 $n=10$ 的观测值的样本：

\begin{mvtable}
\mvfit{%
\begin{tabular}{|c|c|c|c|c|c|}
\hline
顺序观测值 $x_{i(j)}$ & 概率水平 $\left(j-\dfrac{1}{2}\right)\big/n$ & 标准正态分位数 $q_{i(j)}$ & 顺序观测值 $x_{i(j)}$ & 概率水平 $\left(j-\dfrac{1}{2}\right)\big/n$ & 标准正态分位数 $q_{i(j)}$\\
\hline
$-1.00$ & $0.05$ & $-1.645$ & $0.80$ & $0.55$ & $0.125$\\
$-0.10$ & $0.15$ & $-1.036$ & $1.26$ & $0.65$ & $0.385$\\
$0.16$ & $0.25$ & $-0.674$ & $1.54$ & $0.75$ & $0.674$\\
$0.41$ & $0.35$ & $-0.385$ & $1.71$ & $0.85$ & $1.036$\\
$0.62$ & $0.45$ & $-0.125$ & $2.30$ & $0.95$ & $1.645$\\
\hline
\end{tabular}
}
\end{mvtable}
